\documentclass[11pt,letterpaper]{article}
\usepackage{fontspec}
\setmainfont{texgyrepagella}[Extension=.otf,UprightFont=*-regular,BoldFont=*-bold,ItalicFont=*-italic,BoldItalicFont=*-bolditalic]
\usepackage[bidi=default]{babel}
\babelprovide[import,main]{spanish}
\babelprovide[import]{hebrew}
\babelprovide[import]{ancientgreek}
\babelfont[hebrew]{rm}[Path=fonts/,BoldFont=NotoHebrewSerif-700.ttf,Scale=0.95]{NotoHebrewSerif-400.ttf}
\babelfont[ancientgreek]{rm}[Path=fonts/]{GFSPorson.otf}
\newcommand{\heb}[1]{\mbox{\foreignlanguage{hebrew}{#1}}}
\newcommand{\grc}[1]{\foreignlanguage{ancientgreek}{#1}}
\usepackage{needspace,wrapfig}
\newcommand{\hebverso}[1]{\par\needspace{7\baselineskip}\smallskip{\centering\begin{otherlanguage}{hebrew}\large #1\end{otherlanguage}\par}\smallskip}
\usepackage[margin=2.5cm,headheight=14pt]{geometry}
\usepackage{microtype,ragged2e,booktabs,graphicx,xcolor,csquotes,array}
\usepackage[hidelinks]{hyperref}\usepackage{xurl}\urlstyle{same}
\usepackage[most]{tcolorbox}
\usepackage{tikz}\usetikzlibrary{arrows.meta,positioning,calc,decorations.pathreplacing}
\usepackage{titlesec,fancyhdr,extramarks}
% Paleta "Israel": celeste y azul del tejelet (Nm 15:38) con oro del santuario
% (nombres conservados por compatibilidad: ladrillo = azul tejelet, cafe = azul real, cafeosc = azul noche,
%  piedra = celeste claro, sepia = azul pizarra)
\definecolor{cafe}{RGB}{26,82,150}\definecolor{cafeosc}{RGB}{11,33,74}\definecolor{tinta}{RGB}{20,32,56}
\definecolor{oro}{RGB}{206,162,52}\definecolor{ladrillo}{RGB}{20,68,140}
\definecolor{piedra}{RGB}{234,243,251}\definecolor{sepia}{RGB}{72,102,140}\definecolor{celeste}{RGB}{96,166,222}
\titleformat{\section}{\Large\bfseries\color{ladrillo}}{\thesection.}{0.6em}{}
\titleformat{\subsection}{\large\bfseries\color{cafe}}{\thesubsection}{0.6em}{}
\pagestyle{fancy}\fancyhf{}
\fancyhead[L]{\foreignlanguage{spanish}{\small\itshape\color{sepia}\firstleftmark}}\fancyhead[R]{\foreignlanguage{spanish}{\small\color{sepia}\thepage}}
\renewcommand{\sectionmark}[1]{\markboth{#1}{}}
\setlength{\parindent}{1.27cm}\linespread{1.15}\setlength{\emergencystretch}{3em}
\newtcolorbox{cita}[1][]{enhanced,colback=piedra,colframe=oro!90!ladrillo,boxrule=0pt,leftrule=2.5pt,arc=0pt,
  left=8pt,right=8pt,top=5pt,bottom=5pt,before skip=8pt,after skip=8pt,fontupper=\itshape,#1}
\newcommand{\fuente}[1]{\par\hfill{\upshape\bfseries\small\textcolor{ladrillo}{#1}}}
\newtcolorbox{nota}[1][]{enhanced,colback=oro!12,colframe=oro!80!ladrillo,boxrule=0.6pt,arc=2pt,
  fonttitle=\bfseries\small,coltitle=cafe,colbacktitle=oro!35,fontupper=\small,title={#1}}
\newtcolorbox{perla}[1]{enhanced,breakable,colback=oro!7,colframe=oro!85!ladrillo,boxrule=0.8pt,arc=3pt,left=8pt,right=8pt,top=6pt,bottom=6pt,fonttitle=\bfseries\small,coltitle=cafeosc,colbacktitle=oro!32,title={Perla filológica · #1},fontupper=\small,before skip=10pt,after skip=10pt}
\newcommand{\refapa}[1]{\par{\RaggedRight\hangindent=1.27cm\hangafter=1\noindent #1\par}\vspace{4pt}}
\newcommand{\destacar}[1]{\textbf{\textcolor{ladrillo}{#1}}}
